
\documentclass[]{scrreprt}
\usepackage{arbeitsblatt}

\usepackage[bottom=2cm, top=1.5cm]{geometry}
 
\title{Übungsblatt \LaTeX}
\ohead{}
\ihead{Gitkurs}


\author{}
\date{}

\begin{document}


\section*{Erste Schritte im Terminal}
Sicheres Navigieren im Dateisystem sowie Erstellen, Anzeigen und Löschen von
Dateien und Ordnern im Terminal.

\section*{Hinweise}
\begin{itemize}
  \item Erstelle im Explorer ein Ordner mit Namen \texttt{gitkurs}.
  \item Öffne die git-Bash in diesem Ordner (rechte Maus Taste.
  \item Arbeite ausschließlich in \textbf{diesem} Ordner oder tiefer.
  \item Nutze die \texttt{Tab}-Taste zur automatischen Vervollständigung von Befehlen.
  \item Starte die Befehlsausführung immer mit der \texttt{enter}-Taste.
  \item Lies auftretende Fehlermeldungen aufmerksam.
\end{itemize}

\section*{1. Wo bin ich?}

\textbf{Neuer Befehl:}
\begin{lstlisting}
pwd   # kurz für "print working directory"
\end{lstlisting}

\begin{enumerate}
  \item Geben Sie den Befehl ein.
  \item Notiere, welches Verzeichnis angezeigt wird.
\end{enumerate}

\section*{2. Inhalte anzeigen}

\textbf{Neuer Befehl:}
\begin{lstlisting}
ls   # kurz für "list"
\end{lstlisting}

\begin{enumerate}
  \item Lass dir den Inhalt des aktuellen Ordners anzeigen.
  \item Wiederhole den Befehl mit der Option:
\begin{lstlisting}
ls -l    # kurz für "list long"
\end{lstlisting}
  \item Beschreibe den Unterschied.
\end{enumerate}

\section*{3. Verzeichnisse erstellen}

\textbf{Neuer Befehl:}
\begin{lstlisting}
mkdir    # Kurz für "make directory"
\end{lstlisting}

\begin{enumerate}
  \item Erstelle ein Verzeichnis \texttt{shell\_training}.
  \item Wechsle in dieses Verzeichnis.\\ Der Befehl besteht aus den 
  Anfangsbuchstaben von „change directory“.
\end{enumerate}

\section*{4. Dateien erstellen}

\textbf{Neuer Befehl:}
\begin{lstlisting}
touch   # Erstellt leere Datei
\end{lstlisting}

\begin{enumerate}
  \item Erstelle drei leere Dateien:
  \texttt{datei1.txt}, \texttt{datei2.txt}, \texttt{notizen.md}
  \item Kontrolliere das Ergebnis mit einem geeigneten Befehl.
\end{enumerate}

\section*{5. Kopieren und Verschieben}

\textbf{Neue Befehle:}
\begin{lstlisting}
cp  # kurz für "copy"
mv  # kurz für "move". Im gleichen Ordner = umbenennen 
\end{lstlisting}

\begin{enumerate}
  \item Erstelle ein Unterverzeichnis \texttt{backup}.
  \item Kopiere \texttt{datei1.txt} in dieses Verzeichnis.
  \item Verschiebe \texttt{datei2.txt} ebenfalls nach \texttt{backup}.
  \item Prüfe jeweils das Ergebnis.
\end{enumerate}

\section*{6. Löschen}

\textbf{Neue Befehle:}
\begin{lstlisting}
rm     # kurz für "remove"  (Datei)
rmdir  # kurz für "remove directory"
\end{lstlisting}

\begin{enumerate}
  \item Lösche die Datei \texttt{notizen.md}.
  \item Wechsle in das Verzeichnis \texttt{backup}.
  \item Lösche dort eine Datei.
\end{enumerate}

\begin{tcolorbox}[title=Hinweis]
Oben hast du den Befehl \texttt{cd <ordnername>} kennengelernt. Damit kommst du in der Hierarchie aber nur abwärts. Um wieder hinauf zu klettern, kannst du entweder den kompletten Zielpfad angeben oder \texttt{cd ..} verwenden
\end{tcolorbox}



\section*{Abschluss}

Lösche das gesamte Verzeichnis \texttt{shell\_training} inklusive Inhalt.\\
\textbf{Achtung:} Informiere dich vorher über die passende Option von \texttt{rm}.

\end{document}